\section{Como ler um benchmark}
\label{sec:como-ler-benchmark}

\paragraph{O \enquote{melhor de N} engana.}
\textcite{paper:statistically-rigorous-java-perf} partem de uma
observação: o tempo de um programa Java varia de execução para execução,
por causa do JIT (guiado por amostragem de métodos por tempo), do
escalonamento de threads, do GC e do sistema. Mesmo assim, muitos
trabalhos relatavam a média, o melhor ou o segundo melhor de várias
execuções, sem estatística. O exemplo do paper usa a Jikes RVM com cinco
coletores no benchmark \texttt{db} do SPECjvm98: pelo melhor de 30
execuções, CopyMS e GenCopy parecem empatados e SemiSpace parece
claramente melhor que GenCopy; com intervalos de confiança de 95\%,
GenCopy é significativamente melhor que CopyMS, e a diferença entre
SemiSpace e GenCopy pode ser só ruído. O \enquote{melhor de N} relata a execução de
sorte.

\paragraph{Startup e steady state são perguntas diferentes.}
O mesmo paper recomenda medir os dois separadamente:
\begin{description}
\item[Startup] várias medições, cada uma com uma invocação nova da VM e
  uma única iteração do benchmark; intervalo de confiança entre elas.
\item[Steady state] várias invocações da VM, cada uma com várias
  iterações; dentro de cada invocação, só as iterações depois do
  aquecimento; intervalo de confiança entre invocações.
\end{description}
É a mesma divisão da seção~\ref{sec:jvm-vs-nativo}: \emph{time to first
request} é startup, throughput de pico é steady state. Um número sem
dizer qual dos dois mede não diz nada.

\paragraph{Steady state pode não existir.}
\textcite{paper:vm-warmup-hot-cold} testaram a premissa de que, depois do
aquecimento, toda VM chega a um patamar de desempenho de pico. Com
análise de \emph{changepoints}, em três máquinas e em VMs que incluem
HotSpot~8u112b15, Graal~0.18, PyPy, V8 e LuaJIT, encontraram que
\enquote{at most 43.5\% of <VM, Benchmark> pairs consistently reach a
steady state of peak performance}. Alguns benchmarks ficam mais
\emph{lentos} com o tempo. Descartar as primeiras iterações como
\enquote{aquecimento} e tirar a média do resto pode esconder isso.

\paragraph{A carga precisa parecer com a sua.}
O DaCapo nasceu para trocar microbenchmarks por aplicações reais
\parencite{paper:dacapo}; a versão atual, Chopin, pede que se cite a
versão exata usada (\citetitle{doc:dacapo}) e tem paper próprio
\parencite{paper:dacapo-chopin}. O Renaissance
\parencite{paper:renaissance} foi criado porque as suítes anteriores
não exercitavam concorrência e programação funcional modernas (Spark,
Akka, streams), e os autores mostram que otimizações invisíveis no DaCapo
aparecem no Renaissance. Um benchmark de JSON em memória não diz nada
sobre o seu serviço que passa 80\% do tempo esperando o banco.

\paragraph{Checklist para qualquer número.}
\begin{enumerate}
\item \emph{Quem mediu?} Fornecedor, concorrente ou terceiro. Número de
  fornecedor não é falso, mas foi escolhido por quem vende.
\item \emph{Quando, e com que versões?} JDK, framework, VM, sistema.
  Número de 2019 descreve uma VM de 2019.
\item \emph{Em que máquina?} Núcleos, memória, rede, se a aplicação
  dividia a máquina com o gerador de carga.
\item \emph{Mede startup ou steady state?} E, se for steady state, como
  decidiu que chegou lá?
\item \emph{Quantas execuções, e qual a variação?} Sem intervalo de
  confiança ou dispersão, uma diferença de 5\% pode ser ruído.
\item \emph{O que ficou de fora?} Benchmarks excluídos, testes em que o
  sistema perde, configurações ajustadas só de um lado.
\item \emph{A carga parece com a minha?} Se não, use o número para ter
  ordem de grandeza e meça o seu (Parte~I).
\end{enumerate}
